%% ================================================================
\section{Agriculture}
\label{sec:agriculture}

\subsection{Agriculture: scope and decomposition}

Agricultural activities emit a wide range of pollutants through
processes whose spatial drivers differ markedly. Enteric fermentation
and manure management follow livestock systems, synthetic fertiliser
application follows crop mix and management intensity, rice cultivation
is geographically restricted, biomass burning is episodic and
crop-specific, and liming depends on soil acidity rather than land
cover alone. In CAMS-REG these activities are reported primarily under
GNFR~K (agricultural livestock) and GNFR~L (agricultural other), but
their internal heterogeneity is too large to justify a single sectoral
proxy.

For the present thesis, this heterogeneity matters both methodologically
and substantively. Methodologically, agriculture is the first sector in
which the upgraded source-specific framework is implemented in full.
Substantively, the sector contributes to pollutants that are directly
relevant for the later applications, especially NH$_3$ as a precursor
of secondary particulate matter, but also CH$_4$, PM, NMVOC, and
selected CO$_2$ sources that are required for inventory completeness.
The agricultural block is therefore organised into four steps:
identifying the process sources retained, assigning pollutant-specific
mass fractions, defining thematic proxy families, and synthesising the
main uncertainties.

\subsection{Process sources and pollutant mapping}

Table~\ref{tab:pollutant_sources} lists the process sources retained
for each pollutant, together with the GNFR sector to which CAMS assigns
them and the reference used to justify that attribution.

\begin{table}[htbp]
  \centering
  \caption{Agricultural process sources by pollutant and GNFR sector.}
  \label{tab:pollutant_sources}
  \begin{tabular}{lcll}
    \hline
    Pollutant & Source & CAMS sector & Reference \\
    \hline
    CH$_4$ & Enteric fermentation & K & \citep{ref_ch4_nh3_review} \\
     & Manure management & K & \citep{ref_ch4_nh3_review} \\
     & Rice paddies & L & \citep{ref_ch4_nh3_review} \\
    \hline
    NH$_3$ & Synthetic fertiliser application & L & \citep{EMEP_EEA} \\
     & Manure management & K & \citep{EMEP_EEA} \\
     & Livestock housing & K & \citep{EMEP_EEA} \\
    \hline
    NO$_x$ & Fertilised soil & L & \citep{EMEP_EEA} (Chapter~3.D) \\
     & Manure management & K & \citep{EMEP_EEA} (Chapter~3.B) \\
     & Biomass burning & L & \citep{EMEP_EEA} (Chapter~4.F) \\
    \hline
    CO & Biomass burning & L & \citep{andreae_merlet_2001} \\
    \hline
    CO$_2$ & Urea application & L & \citep{IPPC_2006_V10,kim2017} \\
     & Soil liming & L & \citep{IPPC_2006_V10,RB209} \\
    \hline
    NMVOC & Crop production & L & \citep{konig1995,lamb1993} \\
     & Manure management & K & \citep{schade_crutzen_1995,Puliafito_2019} \\
     & Biomass burning & L & \citep{Puliafito_2019} \\
    \hline
    PM$_{10}$, PM$_{2.5}$ & Biomass burning & L & \citep{Puliafito_2019,karanasiou_2021} \\
     & Livestock housing & K & \citep{cambra_lopez_2010} \\
    \hline
  \end{tabular}
\end{table}

The decomposition follows three criteria. First, each source must be
compatible with the way CAMS aggregates emissions at GNFR level.
Second, it must have a spatial behaviour that is distinguishable from
the other sources retained in the same sector. Third, it must matter
for at least one pollutant relevant to the downscaling objective. The
purpose of the decomposition is therefore not to reproduce the full
inventory compilation chain, but to replace an overly generic sectoral
proxy with a set of process-specific spatial patterns that better match
how emissions are expected to occur.

Several agricultural activities are deliberately excluded from this
block because CAMS allocates them to other GNFR sectors. Agricultural
machinery is mapped to GNFR~I (off-road mobile sources) and pesticide
use to GNFR~E (solvent and product use). SO$_2$ emissions from the
agricultural sector are negligible at European scale
(Table~\ref{tab:cams_kl_totals}) and are therefore not decomposed into
sub-sources. CAMS also reports primary particulate matter only, so
secondary organic aerosol formation is not spatially allocated here.

\subsection{Mass-fraction coefficients}

Table~\ref{tab:alpha_agriculture} presents the $\alpha$ values used to
recombine the source-specific proxies into pollutant-specific weights.
For CH$_4$, CO$_2$, CO, and NMVOC, the fractions are derived from
European totals assembled from the EEA reporting framework for 2020
\citep{EEA_GHG_Emissions_2020}. For NH$_3$, NO$_x$, and particulate
matter, the fractions follow the EMEP/EEA Guidebook
\citep{EMEP_EEA}. The one-year mismatch with the 2019 CAMS base year is
accepted here because the coefficients are used as first-order spatial
weights rather than as year-specific source apportionment for Greece.

\begin{table}[htbp]
  \centering
  \caption{European agricultural source totals (Gg) and corresponding
  $\alpha$ values used in the agriculture block.}
  \label{tab:alpha_agriculture}
  \begin{tabular}{lcll}
    \hline
    Pollutant & Source & Emissions & $\alpha$ \\
    \hline
    CH$_4$ & Enteric fermentation & 6575.21 & 0.790 \\
     & Manure management & 1645.17 & 0.198 \\
     & Rice paddies & 98.32 & 0.012 \\
    \hline
    NH$_3$ & Synthetic fertiliser application & 1749 & 0.541 \\
     & Manure-management pathways & 1481 & 0.459 \\
    \hline
    NO$_x$ & Fertilised soil & 638 & 0.936 \\
     & Manure management & 37.6 & 0.055 \\
     & Biomass burning & 5.8 & 0.009 \\
    \hline
    CO$_2$ & Urea application & 4704.57 & 0.468 \\
     & Soil liming & 5346.93 & 0.532 \\
    \hline
    NMVOC & Crop production & 330.96 & 0.239 \\
     & Manure management & 1052.35 & 0.760 \\
     & Biomass burning & 2.04 & 0.001 \\
    \hline
    PM$_{10}$ & Biomass burning & 3.7 & 0.036 \\
     & Livestock housing & 99 & 0.964 \\
    \hline
    PM$_{2.5}$ & Biomass burning & 3.6 & 0.146 \\
     & Livestock housing & 21 & 0.854 \\
    \hline
  \end{tabular}
\end{table}

The combined NH$_3$ row is retained deliberately. In the EMEP
accounting logic, manure storage and application are treated together,
and the total manure-management contribution also includes manure
generated in open yards and excreta deposited during grazing. A more
granular split between manure management and housing would therefore
require assumptions that are not directly supported by the European
totals used here. The combined coefficient is consequently interpreted
as a pragmatic budget split, while the spatial allocation itself is
still improved by separating the corresponding proxy patterns.

For carbon monoxide, biomass burning is the only agricultural source
retained in this framework, so no source competition arises within the
sectoral budget. The agriculture subsections below therefore focus on
sources that differ spatially rather than artificially subdividing a
single-source pollutant.

\subsection{Livestock-related proxies}

The first proxy family targets emissions controlled by the geography of
animal production systems. These sources are tied to grazing, herd
composition, manure handling, and confined livestock infrastructure.

\subsubsection{Enteric fermentation}
\label{sec:mu-rho-enteric}

\paragraph{Rationale and spatial driver.}
Enteric CH$_4$ is produced in the rumen and hindgut of ruminants. Its
spatial pattern is driven first by where animals graze and second by
which ruminant species dominate in a given region. The proxy therefore
combines direct field evidence of grazing from LUCAS with species-level
livestock counts from the Eurostat agricultural census.

\paragraph{Input data and eligibility.}
The point-level signal comes from the LUCAS grazing indicator and the
LUCAS land-cover code \texttt{LC1}; the regional intensity modifier
comes from Eurostat Agricultural Census livestock counts at NUTS-2
level (Appendix~\ref{app:c21}). Only points falling in agricultural CLC
classes are retained.

\paragraph{Score definition.}
Each LUCAS point receives a grazing score
$s_i^{\mathrm{ent}}$, where $\gamma$ denotes the
\texttt{SURVEY\_GRAZING} code:
\begin{equation}
  s_i^{\mathrm{ent}} =
  \begin{cases}
    1.0, & \gamma = 1 \text{ (grazing observed)},\\[3pt]
    0.0, & \gamma \in \{0,2\} \text{ (no grazing observed)},\\[3pt]
    0.6, & \gamma \text{ missing and } \texttt{LC1} \in \{\text{E*},\text{B55}\},\\[3pt]
    \text{excluded}, & \text{otherwise}.
  \end{cases}
\end{equation}
The fallback value of 0.6 is used as a moderate expert prior for
grassland points without a grazing record: it preserves a positive
signal where grazing is agronomically plausible, while keeping the
score clearly below a confirmed observation.

For each NUTS-2 region $n$ and CLC class $c$, the mean grazing score is
\begin{equation}
  \bar{s}_{n,c}^{\mathrm{ent}} =
  \frac{1}{|\mathcal{P}_{n,c}|}
  \sum_{i \in \mathcal{P}_{n,c}} s_i^{\mathrm{ent}},
\end{equation}
where $\mathcal{P}_{n,c}$ is the set of retained points.

To account for differences in ruminant composition, the mean score is
scaled by a livestock-intensity index:
\begin{equation}
  I_n^{\mathrm{rum}} =
  116\,N_{\mathrm{bovine}} + 8\,N_{\mathrm{sheep}} + 5\,N_{\mathrm{goats}},
\end{equation}
where the coefficients are IPCC~(2006) default enteric CH$_4$
emission factors in kg~CH$_4$~head$^{-1}$~yr$^{-1}$
\citep{IPPC_2006_V10}. The corresponding weight is
\begin{equation}
  \omega_n =
  \mathrm{clip}\!\left(
    \frac{I_n^{\mathrm{rum}}}{\tilde{I}_{\mathrm{EU}}},
    0.05, 5.0
  \right).
\end{equation}
The clipping range regularises the proxy so that a few very high-density
regions do not dominate the full European scaling.

\paragraph{Aggregation and normalisation.}
The raw proxy is
\begin{equation}
  \mu_{n,c}^{\mathrm{ent}} =
  \omega_n \cdot \bar{s}_{n,c}^{\mathrm{ent}},
\end{equation}
and the corresponding relevance field follows
Equation~\ref{eq:rho-def}. This subsection captures \emph{animal
presence}; the next one addresses where the associated excreta are
stored, handled, or applied.

\paragraph{Limitations and expected bias.}
The LUCAS grazing indicator does not distinguish between ruminant
species, so species-specific information enters only through the NUTS-2
scaling stage and cannot resolve within-region herd composition. The
fallback score of 0.6 remains an agronomic prior rather than a direct
observation, and the clipping thresholds are regularisation parameters
rather than measured physical limits.

\subsubsection{Manure management}
\label{sec:mu-rho-manure}

\paragraph{Rationale and spatial driver.}
Manure-management emissions (CH$_4$, NH$_3$, NO$_x$, NMVOC) arise from
storage, handling, land application, and direct deposition of excreta
on pasture. Their spatial variation therefore depends both on where
livestock are present and on which agricultural parcels are likely to
receive manure.

\paragraph{Input data and eligibility.}
The proxy combines two signals from LUCAS: the grazing score already
used for enteric fermentation, which captures direct pasture
deposition, and a land-use weight based on \texttt{SURVEY\_LU1} and
\texttt{SURVEY\_LC1}, which approximates the likelihood of manure
application or livestock-related manure handling on the parcel.

\paragraph{Score definition.}
The grazing component is identical to the enteric score
$s_i^{\mathrm{ent}}$. In the manure context it captures direct
deposition on pasture, which \citet{oenema_2017} identify as the single
largest manure-management category in European livestock systems,
accounting for roughly 30--40\,\% of total manure-management activity.
The EMEP/EEA Guidebook \citep{EMEP_EEA} further notes that NH$_3$
emission factors per kg\,N excreted are substantially lower for pasture
deposition than for housing, storage, and land-application pathways,
because urine and dung deposited on grass volatilise more slowly than
slurry handled in confined systems. The grazing component therefore
carries a stronger signal for CH$_4$ than for NH$_3$, but it remains
necessary for the manure proxy because direct deposition is part of the
inventory source structure.

The second component represents likely land application. This is not
observed directly in LUCAS, so the weighting scheme is necessarily
inferential. \citet{Velthof2009} allocate ruminant manure to fodder
land at approximately equal per-hectare rates and distinguish different
application rates across arable crop groups for monogastric manure.
That pattern supports assigning relatively high weights to managed
grassland and to arable classes associated with crops that commonly
receive manure, while keeping all weights below unity because the proxy
expresses likelihood rather than confirmed spreading. The land-use
weights used for the second component are:
\begin{itemize}
  \setlength{\itemsep}{0.25em}
  \setlength{\parskip}{0pt}
  \item managed grassland in agricultural use: $w_{\mathrm{LU}}=0.8$;
  \item high manure-use arable classes
        $(\texttt{B11}, \texttt{B13}, \texttt{B2*}, \texttt{B32})$:
        $w_{\mathrm{LU}}=0.8$;
  \item intermediate manure-use arable classes
        $(\texttt{B12}, \texttt{B14}, \texttt{B15}, \texttt{B16},
        \texttt{B17}, \texttt{B18}, \texttt{B19}, \texttt{B31})$:
        $w_{\mathrm{LU}}=0.6$;
  \item fallow land: $w_{\mathrm{LU}}=0.7$;
  \item other agricultural land with only limited or uncertain manure
        reception: $w_{\mathrm{LU}}=0.1$.
\end{itemize}
Managed grassland receives a high weight because fodder land is a
primary destination for ruminant manure in the European nitrogen-budget
framework of \citet{Velthof2009}. Within cropland, the high and
intermediate groups follow the crop differentiation in that same study,
which separates classes with strong evidence of manure use from those
where application is less common. Fallow land is kept relatively high
because manure can be applied before the next cropping cycle, whereas
permanent crops and other ambiguous agricultural land remain low because
their manure reception is more variable and less consistently documented
\citep{oenema_2017}. The resulting weights are therefore
literature-informed heuristics interpreted as probabilities of likely
manure reception rather than direct observations of application.

The point-level manure score is then
\begin{equation}
  s_i^{\mathrm{man}} =
  \begin{cases}
    \max\!\left(s_i^{\mathrm{ent}}, w_{\mathrm{LU}}\right),
      & \text{if } s_i^{\mathrm{ent}} \text{ is defined},\\[4pt]
    w_{\mathrm{LU}},
      & \text{if } s_i^{\mathrm{ent}} \text{ is undefined and }
        w_{\mathrm{LU}} > 0,\\[4pt]
    \text{excluded}, & \text{otherwise}.
  \end{cases}
\end{equation}
The maximum is preferred to a sum because direct deposition and
land-application evidence are not independent pathways at point scale;
adding them would double-count precisely those parcels where the
livestock signal is already strongest.

\paragraph{Aggregation and normalisation.}
The raw proxy is the arithmetic mean of $s_i^{\mathrm{man}}$ over
retained points in each pair $(n,c)$:
\begin{equation}
  \mu_{n,c}^{\mathrm{man}} =
  \frac{1}{|\mathcal{P}_{n,c}|}
  \sum_{i \in \mathcal{P}_{n,c}} s_i^{\mathrm{man}},
\end{equation}
with $\mu_{n,c}^{\mathrm{man}}=0$ for empty pairs. Relevance then
follows Equation~\ref{eq:rho-def}. This proxy is intended for spatial
allocation only, not for process-level manure accounting.

\paragraph{Limitations and expected bias.}
The score is pollutant-agnostic even though CH$_4$, NH$_3$, NO$_x$, and
NMVOC respond differently to pasture deposition, slurry, storage type,
and application technique. The resulting proxy is therefore a
deliberate compromise across the four pollutants. National manure
statistics or farm-register data would allow a more differentiated
treatment. The next subsection narrows the focus to emissions
associated with confined livestock infrastructure.

\subsubsection{Livestock housing (NH$_3$ and primary PM)}
\label{sec:proxy-housing}

\paragraph{Rationale and spatial driver.}
This proxy represents confined livestock housing and associated yard or
storage activity, which is used here for NH$_3$ and for housing-related
primary PM. The same spatial pattern is considered acceptable for both
pollutant groups because they are colocated around barns, animal yards,
and manure-handling surfaces, while pollutant-specific differences are
handled through the $\alpha$ coefficients rather than through separate
geographies. In the NH$_3$ case, the process corresponds to confined
livestock housing and associated manure storage within GNFR~K, excluding
excreta deposited directly during grazing. Emissions arise when urine,
faeces, and soiled bedding accumulate in enclosed or semi-enclosed
systems such as barns, animal yards, and slurry stores, where
volatilisation is enhanced by ventilation, temperature, and exposed
emitting surfaces \citep{EMEP_EEA,wyer_2022}. Spatial variation is
therefore controlled by the geography of livestock production, the
degree of animal confinement, and the contrast between grazing-based and
housed systems. In LUCAS~2022, agricultural activity is represented
primarily through \texttt{SURVEY\_LU1}, with \texttt{U111} indicating
agricultural production, while \texttt{SURVEY\_LC1} helps distinguish
likely livestock settings from crop-dominated parcels.

\paragraph{Input data and eligibility.}
The proxy uses \texttt{SURVEY\_GRAZING}, \texttt{SURVEY\_LU1}, and
\texttt{SURVEY\_LC1} from LUCAS, together with gridded livestock counts
for cattle, pigs, poultry, sheep, and goats. Points are retained when
they indicate an agricultural setting plausibly associated with housed
livestock activity and when they do not simply represent direct grazing
deposition.

\paragraph{Score definition.}
The stage-1 point score applies a grazing filter followed by a
land-use/land-cover tier. Let $\gamma$ again denote the
\texttt{SURVEY\_GRAZING} code. If grazing is directly observed
($\gamma=1$), then $s_i^{\mathrm{hous}}=0$ because the point is
interpreted as an open grazing context. If grazing is not observed
($\gamma=2$), the score takes one of the following values:
\begin{itemize}
  \setlength{\itemsep}{0.2em}
  \setlength{\parskip}{0pt}
  \item farm-building context within agriculture: 1.00;
  \item managed grassland in agricultural use: 0.55;
  \item cropland in agricultural use: 0.20;
  \item ley pasture or temporary grass within cropland: 0.30;
  \item fallow land: 0.20.
\end{itemize}
If grazing is ambiguous or unrecorded, the same tier is discounted by
$\lambda=0.4$:
\begin{equation}
  s_i^{\mathrm{hous}} = \lambda \cdot s_i^{\mathrm{tier}}.
\end{equation}
The factor $\lambda$ is retained as a calibrated heuristic to keep
ambiguous points below otherwise equivalent non-grazing observations.
Setting $\lambda=0.4$ also keeps the treatment consistent with the
broader logic used elsewhere in the chapter: when no direct grazing
signal is observed, the housing score is damped rather than interpreted
as positive evidence of animal presence. A non-zero cropland weight is
still retained because some livestock systems in LUCAS are embedded in
mixed crop--livestock settings, where cropland parcels may sit close to
yards, housing, or manure storage even if the parcel itself is not a
building footprint; the weight is therefore reduced to 0.2 rather than
set to zero.

The stage-1 signal is aggregated as
\begin{equation}
  \bar{s}_{n,c}^{\mathrm{hous}} =
  \frac{1}{|\mathcal{P}^{*}_{n,c}|}
  \sum_{i \in \mathcal{P}^{*}_{n,c}} s_i^{\mathrm{hous}}.
\end{equation}

The stage-2 livestock-intensity factor is
\begin{equation}
  I_g^{\mathrm{hous}} =
  41.16\,N_{\mathrm{bovine},g}
  + 8.00\,N_{\mathrm{pigs},g}
  + 0.54\,N_{\mathrm{poultry},g}
  + 0.291\,N_{\mathrm{sheep},g}
  + 0.231\,N_{\mathrm{goats},g},
\end{equation}
where the coefficients are taken from
Table~\ref{tab:livestock_housing_factors}. The corresponding weight is
\begin{equation}
  \omega_g =
  \mathrm{clip}\!\left(
    \frac{I_g^{\mathrm{hous}}}{\tilde{I}_{\mathrm{EU}}},
    0.05, 5.0
  \right).
\end{equation}
As in the enteric proxy, the clipping range acts as a regularisation
device that prevents a small number of intense livestock cells from
dominating the proxy.

\paragraph{Aggregation and normalisation.}
The raw proxy is
\begin{equation}
  \mu_{n,c}^{\mathrm{housing}} =
  \omega_{g(n,c)} \cdot \bar{s}_{n,c}^{\mathrm{hous}},
\end{equation}
where $g(n,c)$ denotes the ancillary livestock grid cell used for the
pair $(n,c)$. Relevance follows Equation~\ref{eq:rho-def}. The proxy is
thus complementary to the grazing-based pathways rather than
overlapping with them: observed grazing suppresses housing relevance,
whereas agricultural points without grazing evidence retain positive
relevance according to how strongly their context suggests confined
livestock activity.

\paragraph{Limitations and expected bias.}
The proxy does not distinguish housing technologies, storage systems,
or production regimes such as dairy versus beef or intensive versus
extensive husbandry. It also infers housing from land-use context
rather than directly observing buildings or animal capacity. The next
proxy family therefore turns away from livestock infrastructure and
towards crop and fertiliser management.

\subsection{Crop and fertiliser-related proxies}

The second family covers sources whose spatial pattern follows crop
composition, nutrient-management intensity, and fertiliser form rather
than livestock location alone. Two distinct fertiliser pathways are
retained on purpose. The first proxy targets the general geography of
mineral-nitrogen use and is used for NH$_3$ and NO$_x$, whose spatial
variation is controlled primarily by where mineral fertiliser is applied
and how intensively it is used. The second targets urea specifically
and is retained because CO$_2$ from urea hydrolysis is a product-
specific source that cannot be represented by a general mineral-N
intensity field alone. The two subsections should therefore be read as
complementary rather than redundant: one is a broad management-intensity
proxy, the other a product-specific proxy.

\subsubsection{Fertilised soil (synthetic nitrogen)}
\label{sec:mu-rho-synth}

\paragraph{Rationale and spatial driver.}
NH$_3$ and NO$_x$ associated with mineral nitrogen application vary
primarily with crop type, fertiliser rate, and management intensity.
The present proxy does not estimate emissions directly; instead, it
uses literature-derived mineral-N application rates as a spatial
indicator of relative emission intensity. That makes the proxy
appropriate for allocation of a known sectoral total even when absolute
rates are uncertain. Its role in the chapter is to represent the
\emph{general} geography of mineral-fertiliser use, irrespective of the
exact commercial product applied. The underlying chemistry differs by
pathway.
Ammoniacal fertilisers and the ammonium formed after urea hydrolysis can
release NH$_3$ through volatilisation:
\begin{equation}
  \mathrm{NH_4^+ + OH^- \rightleftharpoons NH_3 + H_2O}.
\end{equation}
In soils, nitrification and related microbial transformations of
ammonium also generate NO and, more broadly, NO$_x$:
\begin{equation}
  \mathrm{NH_4^+ + 2\,O_2 \rightarrow NO_3^- + 2\,H^+ + H_2O},
\end{equation}
\begin{equation}
  \mathrm{NO \xrightarrow[\text{atmosphere}]{} NO_2},
\end{equation}
while denitrification can also yield NO under favourable soil
conditions. The proxy therefore represents the spatial intensity of
mineral-N use as the common driver of these nitrogen-loss pathways,
rather than trying to reconstruct each reaction mechanism separately at
plot scale.

\paragraph{Input data and eligibility.}
Eligible LUCAS points are those whose \texttt{LC1} indicates cropland
(\texttt{B*}), permanent grassland (\texttt{E10}), or semi-natural
grassland (\texttt{E20}) combined with agricultural land use
\texttt{U111}, which is used here as a proxy for improved grassland
receiving mineral nitrogen. Points whose LC1 code cannot be mapped to
the rate table are excluded.

\paragraph{Score definition.}
Each eligible point is mapped to a crop category $k$ and assigned the
literature rate $R_k$ from Table~\ref{tab:synthetic-n-rates}:
\begin{equation}
  s_i^{\mathrm{synth}} = R_{k(\texttt{LC1},\texttt{LU1})}.
\end{equation}
The rates are expressed in kg~N~ha$^{-1}$ and are not bounded to
$[0,1]$. This is deliberate: the subsequent normalisation operates on
relative ranking rather than on absolute magnitude.

\paragraph{Aggregation and normalisation.}
The raw proxy is
\begin{equation}
  \mu_{n,c}^{\mathrm{synth}} =
  \frac{1}{|\mathcal{P}_{n,c}|}
  \sum_{i \in \mathcal{P}_{n,c}} s_i^{\mathrm{synth}},
\end{equation}
with $\mu_{n,c}^{\mathrm{synth}}=0$ when no eligible points are
available. Relevance then follows Equation~\ref{eq:rho-def}. The next
subsection uses this broad synthetic-N logic as a starting point, but
restricts the interpretation to the specific case of urea. Keeping both
subsections is methodologically justified because they represent two
different pathways: the present proxy captures the general geography of
mineral-N application relevant for NH$_3$ and NO$_x$, whereas the urea
proxy isolates the product-specific CO$_2$ source that arises only when
the applied fertiliser contains urea. In other words, the synthetic-N
proxy answers the question ``where is mineral fertilisation likely to be
intense?'', whereas the urea proxy answers the narrower question
``where is that mineral fertilisation more likely to involve urea?''.

\paragraph{Limitations and expected bias.}
The rates are European literature defaults rather than year-specific,
farm-specific observations, and the LC1-to-category mapping may miss
rare crop classes. The proxy should therefore be interpreted as a
relative fertilisation-intensity field, not as a reconstructed
fertiliser budget.

\subsubsection{Urea application}
\label{sec:mu-rho-urea}

\paragraph{Rationale and spatial driver.}
This subsection represents the spatial distribution of CO$_2$ emissions
from \textbf{urea application to agricultural soils}. Unlike
Section~\ref{sec:mu-rho-synth}, which represents the overall intensity
of \emph{mineral nitrogen} use across crop and grassland systems, the
present proxy focuses on the subset of mineral fertiliser applied in
the form of urea. This distinction matters because the relevant
emission mechanism is specific to urea itself rather than to synthetic
nitrogen in general.

Following the IPCC methodology, urea hydrolysis ultimately releases
carbon dioxide:
\begin{equation}
  \mathrm{CO(NH_2)_2 + H_2O \rightarrow 2\,NH_3 + CO_2}.
\end{equation}
Accordingly,
\begin{equation}
  E_{\mathrm{CO_2},\mathrm{urea}} =
  M_{\mathrm{urea}} \cdot EF_{\mathrm{urea}},
\end{equation}
where $M_{\mathrm{urea}}$ is the amount of urea applied and
$EF_{\mathrm{urea}}$ is the IPCC emission factor
\citep{IPPC_2006_V10}. Because harmonised subnational statistics on
urea use are unavailable, the proxy is framed explicitly as a
first-order approximation of likely urea-use geography rather than as a
direct observation of product-specific application. Spatial variation in
this source therefore depends not simply on where agricultural land
receives mineral nitrogen, but on where urea specifically is likely to
be used. In practice, the proxy is constructed from two elements:
first, a baseline distinction between cropland and grassland from LUCAS
land-cover observations; and second, a grassland-specific livestock-
intensity adjustment intended to capture differences in the management
intensity of fodder and grazing systems. This subsection therefore does
not replace the previous synthetic-N proxy. Instead, it inherits the
idea that fertiliser use intensity is spatially structured, but narrows
the question to the subset of that activity involving urea, because
only that subset produces the CO$_2$ source of interest here.

\paragraph{Input data and eligibility.}
Retained LUCAS points are classified as cropland or grassland
according to the agricultural land-cover selection described in
Appendix~\ref{app:lucas}. Gridded bovine, sheep, and goat counts are
used only for the grassland adjustment, where management intensity is
expected to vary most strongly. This separation mirrors the conceptual
structure of the proxy: cropland provides the baseline domain in which
synthetic fertiliser use is most common, while grassland requires an
additional intensity term because fertiliser use there is more
heterogeneous across extensive and intensive systems.

\paragraph{Score definition.}
The baseline point score is
\begin{equation}
  s_{i,0}^{\mathrm{urea}} =
  \begin{cases}
    1.0, & \text{cropland},\\[3pt]
    0.7, & \text{grassland},\\[3pt]
    \text{excluded}, & \text{otherwise}.
  \end{cases}
\end{equation}
The cropland value is set to 1.0 because cropped land is the primary
domain of synthetic fertiliser application. The lower grassland value
does not imply a fixed ratio of urea use between cropland and
grassland; it is a heuristic way of expressing that fertilised
grassland is both eligible and more heterogeneous. This baseline
distinction is intended to capture the fact that mineral fertiliser use
extends to pastures, but less uniformly than on cropland
\citep{IFASTAT_consumption,einarsson_2021}.

For grassland points, the baseline is scaled by a livestock-intensity
indicator. The motivation is to approximate how strongly a given
grassland area is associated with surrounding livestock production
systems. In agricultural regions where more grazing livestock are
supported, grassland and fodder areas are more likely to be managed
intensively and to require higher nutrient inputs, including mineral
nitrogen application. This relationship is not treated as a direct
observation of urea use, but as a proxy for the relative intensity of
managed grassland systems. This interpretation is consistent with
Eurostat's use of grazing livestock density as an indicator defined
relative to fodder area, including forage crops and permanent
grassland. Eurostat further notes that higher grazing livestock density
implies more manure available per hectare of fodder area and therefore
greater nutrient pressure, even though realised impacts remain
management-dependent
\citep{EurostatLivestockPatterns,EurostatGrazingGlossary}. The indicator
is therefore written as:
\begin{equation}
  I_g^{\mathrm{grass}} =
  N_{\mathrm{bovine},g} +
  N_{\mathrm{sheep},g} +
  N_{\mathrm{goats},g},
\end{equation}
\begin{equation}
  \omega_g^{\mathrm{grass}} =
  \frac{I_g^{\mathrm{grass}}}{\tilde{I}_{\mathrm{EU}}^{\mathrm{grass}}}.
\end{equation}
The final point score is
\begin{equation}
  s_i^{\mathrm{urea}} =
  \begin{cases}
    1.0, & i \text{ is cropland},\\[3pt]
    0.7 \cdot \omega_{g(i)}^{\mathrm{grass}}, & i \text{ is grassland},\\[3pt]
    \text{excluded}, & \text{otherwise}.
  \end{cases}
\end{equation}
Livestock intensity is applied only to grassland because the proxy uses
it as an indicator of managed fodder-system intensity, not as a direct
determinant of urea use everywhere. The cropland baseline is therefore
left unchanged, while the grassland branch is allowed to vary according
to the relative intensity of the surrounding grazing-livestock system.

\paragraph{Aggregation and normalisation.}
The raw proxy is
\begin{equation}
  \mu_{n,c}^{\mathrm{urea}} =
  \frac{1}{|\mathcal{P}_{n,c}|}
  \sum_{i \in \mathcal{P}_{n,c}} s_i^{\mathrm{urea}},
\end{equation}
with $\mu_{n,c}^{\mathrm{urea}}=0$ for empty pairs. Relevance then
follows Equation~\ref{eq:rho-def}. Relative to the previous subsection,
the key distinction is that the synthetic-N proxy represents broad
fertilisation intensity, whereas the present proxy isolates only the
subset of that management likely to involve urea. The pairing is thus
hierarchical rather than duplicative: the synthetic-N proxy describes
the broader fertilisation domain, while the urea proxy selects one
product-specific pathway within that domain.

\paragraph{Limitations and expected bias.}
The proxy does not observe urea directly and cannot distinguish it from
other mineral-N products such as ammonium nitrate or
urea-ammonium-nitrate solutions. The grassland scaling is also only an
indicator of management intensity, not a one-to-one predictor of urea
application. Future work could improve this proxy through fertiliser
sales data, farm-accountancy data, or crop-specific information on the
share of urea in total mineral-N use.

\subsubsection{Crop production NMVOC}
\label{sec:proxy-agsoils-nmvoc}

\paragraph{Rationale and spatial driver.}
This proxy represents the crop-production component of the agricultural
NMVOC inventory category. Although the reported category is often named
``crop production and agricultural soils'', the spatial signal retained
here is driven primarily by vegetation type and crop composition rather
than by soil properties. Spatial variation is therefore assumed to
follow crop area and crop-specific emission factors.

The EMEP/EEA Guidebook \cite{EMEP_EEA} expresses emissions as
\begin{equation}
  E_{\mathrm{NMVOC}} = \sum_i A_i \times m_{D,i} \times t_i \times EF_i,
  \label{eq:emep-nmvoc}
\end{equation}
where $A_i$ is crop area, $m_{D,i}$ is mean dry-matter yield, $t_i$ is
emission duration, and $EF_i$ is the crop-specific factor. The proxy
therefore uses crop composition as the dominant control.

\paragraph{Input data and eligibility.}
All agricultural CORINE points (CLC~12--22) with a valid
\texttt{SURVEY\_LC1} code are eligible. Points with missing crop
information are excluded. The crop-to-factor mapping follows the EMEP
Guidebook and the Th\"unen NIR rule \citep{thuenen_nmvoc}, under which
only a small number of crop groups have directly measured factors and
other crop types inherit either the wheat or grass factor. This
assignment is not arbitrary: the Guidebook Table~3.3 and Annex~A3
ultimately trace back to measurements reported by \citet{konig1995} and
\citet{lamb1993}, while the Th\"unen methodology provides an explicit
inventory rule for extending those sparse measurements to broader crop
classes in a way that remains consistent with operational national
inventory practice.

\paragraph{Score definition.}
For each LUCAS point $i$ with LC1 code $k$,
\begin{equation}
  s_i^{\mathrm{nmvoc}} =
  EF_k \times \frac{\texttt{SURVEY\_LC1\_PERC}}{100},
  \label{eq:nmvoc-point}
\end{equation}
where $EF_k$ comes from Table~\ref{tab:nmvoc-ef}. If LC1 has no entry
in the table, the point receives zero. A secondary check on
\texttt{SURVEY\_LC2} is applied if LC1 is not matched. The justification
for assigning wheat as the default EF to many annual crops is that the
available literature does not support crop-specific factors for the full
European crop spectrum; the Th\"unen rule is therefore used as a
transparent and reproducible bridge between sparse measurements and the
broader crop nomenclature required by the proxy.

\paragraph{Aggregation and normalisation.}
The raw proxy is
\begin{equation}
  \mu_{n,c}^{\mathrm{nmvoc}} =
  \frac{1}{N_{n,c}} \sum_{i \in (n,c)} s_i^{\mathrm{nmvoc}},
\end{equation}
with $\mu_{n,c}^{\mathrm{nmvoc}}=0$ for empty pairs. Unlike biomass
burning, no external spatial dataset is required because the crop
composition observed by LUCAS is itself the relevant spatial signal.

\paragraph{Limitations and expected bias.}
Only four crop groups have independently measured EFs, and the
remaining categories inherit proxy factors. This introduces substantial
uncertainty, especially for permanent crops in southern Europe, which
are treated as non-emitting in the present implementation and may
therefore be underestimated. The next family turns from crop and
fertiliser management to specialised land-use cases whose geography is
more sharply constrained.

\subsection{Specialised land-use proxies}

The final family contains sources whose spatial pattern is dominated by
highly specific land uses or environmental conditions rather than by
general agricultural management intensity.

\subsubsection{Rice paddies (CH$_4$)}
\label{sec:proxy-rice}

\paragraph{Rationale and spatial driver.}
Methane from flooded rice soils arises from methanogenesis in anoxic,
carbon-rich pore water. Emission factors depend strongly on water
regime, with \citet{yann2009} reporting a difference of roughly a
factor of three between irrigated rice and rain-fed or deepwater
systems. Within Europe, however, the first-order determinant of spatial
variation is simply where rice is grown, because cultivation is
geographically concentrated in a limited number of regions such as the
Po Valley, the Ebro Delta, the Camargue, the Thessaloniki plain, and
the Mondego lowlands \citep{bogliani_2014}. A presence proxy is
therefore a defensible first approximation, even though it does not
resolve the more detailed effect of water management.

\paragraph{Input data and eligibility.}
Only LUCAS points identified as rice by \texttt{LC1}=\texttt{B17} are
retained. All other points are excluded.

\paragraph{Score definition.}
The point-level score is binary:
\begin{equation}
  s_i^{\mathrm{rice}} =
  \begin{cases}
    1.0, & \text{if } \texttt{LC1}=\texttt{B17},\\
    \text{excluded}, & \text{otherwise}.
  \end{cases}
\end{equation}
This is explicitly a presence proxy, not an intensity proxy. Variation
in water management is left to the emission factor embedded in the
inventory.

\paragraph{Aggregation and normalisation.}
The raw proxy is
\begin{equation}
  \mu_{n,c}^{\mathrm{rice}} =
  \frac{1}{|\mathcal{P}_{n,c}|}
  \sum_{i \in \mathcal{P}_{n,c}} s_i^{\mathrm{rice}},
\end{equation}
with zero assigned where no retained points exist. Relevance follows
Equation~\ref{eq:rho-def}. The next subsection remains land-use
specific, but combines remote sensing with crop composition rather than
using field observations alone.

\paragraph{Limitations and expected bias.}
The proxy captures the presence of rice cultivation but not the
intensity of flooding or drainage management. Satellite-derived
inundation products or national rice calendars would be required to
resolve that additional variability.

\subsubsection{Biomass burning (field residues)}
\label{sec:proxy-biomass}

\paragraph{Rationale and spatial driver.}
Open burning of crop residues emits NO$_x$, CO, PM$_{2.5}$, PM$_{10}$,
CH$_4$, and other combustion products directly into the lower
atmosphere. As emphasised by \citet{pinakana2024} and
\citet{andreae_2019}, residue burning is often selected because it is a
rapid and low-cost way to clear fields before the next sowing cycle,
even though it carries clear air-quality and health implications. The
practice is formally prohibited across the EU under GAEC~3
\citep{eu_reg_2021_2115}, yet satellite-derived burned-area products
continue to identify agricultural fire signals in cropland areas across
Europe \citep{forrest_2024}. This combination of legal prohibition and
observed persistence makes farmer surveys or administrative declarations
insufficient as the sole spatial basis. A credible proxy must therefore
identify both where agricultural fires occur and what residues are
available to burn.

\paragraph{Input data and eligibility.}
The proxy combines GFED4.1s satellite fire emissions with LUCAS crop
composition. GFED provides the regional fire geography, while LUCAS
identifies the residue-producing crop mix within agricultural classes.

\paragraph{Score definition.}
The annual agricultural fire score for GFED pixel $r$ is
\begin{equation}
  s_r^{\mathrm{burn}} =
  \frac{1}{N_{\mathrm{yr}}}
  \sum_{\mathrm{yr}=1997}^{2015}
  \sum_{m=1}^{12}
  E_{r,m}^{\mathrm{yr}} \cdot f_{r,m}^{\mathrm{AGRI,yr}} \cdot A_r,
  \label{eq:gfed-score}
\end{equation}
where $A_r$ is the pixel area. The 19-year mean is used to recover a
stable climatological spatial pattern rather than a single anomalous
year and is therefore treated as a proxy for persistent fire-prone
locations in the absence of a harmonised annual European product for
the inventory year.

At NUTS-2 level,
\begin{equation}
  \mu_n^{\mathrm{burn}} =
  \sum_{r:\,\mathrm{centre}(r)\in n} s_r^{\mathrm{burn}}.
  \label{eq:nuts2-burn}
\end{equation}

Within each pair $(n,c)$, residue availability is represented by
\begin{equation}
  \bar{w}(n,c) =
  \frac{1}{N_{n,c}}
  \sum_{k \in \mathcal{C}_{\mathrm{emep}}} N_{n,c,k} \cdot s_k,
  \label{eq:lucas-weight}
\end{equation}
where $\mathcal{C}_{\mathrm{emep}}$ is the eligible crop set and $s_k$
the crop-specific residue-to-yield ratio. The final raw proxy is
\begin{equation}
  \mu_{n,c}^{\mathrm{burn}} =
  \mu_n^{\mathrm{burn}} \cdot
  \frac{n_{n,c}\,\bar{w}(n,c)}
       {\displaystyle\sum_{c'} n_{n,c'}\,\bar{w}(n,c')}.
  \label{eq:clc-split}
\end{equation}

\paragraph{Aggregation and normalisation.}
The corresponding relevance field is obtained from
Equation~\ref{eq:rho-def}. Unlike the preceding proxies, biomass
burning merges a coarse remote-sensing climatology with fine-scale crop
composition, which is why the within-region split is essential.

\paragraph{Limitations and expected bias.}
GFED4.1s ends in 2015, so temporal transferability is assumed when the
climatology is applied to the later CAMS inventory year. The proxy also
inherits uncertainty from crop rotations and from occasional
misclassification of non-agricultural fires near field boundaries. The
last specialised proxy differs again by relying on direct soil
observations rather than crop or fire data.

\subsubsection{Soil liming (CO$_2$)}
\label{sec:proxy-liming}

\paragraph{Rationale and spatial driver.}
Agricultural liming emits CO$_2$ when carbonate amendments are applied
to acidic soils. The dominant spatial control is therefore soil pH
rather than crop area alone. This makes liming a useful example of a
proxy that cannot be represented well by land cover without auxiliary
soil data. NFR~3.G.1 covers the application of limestone
(\(\mathrm{CaCO_3}\)) or dolomite (\(\mathrm{CaMg(CO_3)_2}\)) to acidic
soils, with carbon dioxide released during neutralisation of soil
acidity:
\begin{equation}
  \mathrm{CaCO_3 + 2\,H^+ \rightarrow Ca^{2+} + H_2O + CO_2}.
\end{equation}
The spatial pattern is governed primarily by two factors. First, soil
pH determines whether liming is agronomically required at all: RB209
\citep{RB209} sets crop- and soil-specific target pH values and
therefore provides an operational criterion for when lime demand should
fall to zero. Second, land use matters because arable systems generally
require a higher target pH than managed grassland, while strongly
calcareous soils typically require no lime application
\citep{goulding2016,RB209}.

\paragraph{Input data and eligibility.}
The proxy uses directly measured soil pH from LUCAS Soil 2018
(Appendix~\ref{app:lucas_topsoil}) together with the LUCAS land-cover
code and soil organic carbon. Points are retained if they correspond to
cropland (\texttt{B*}) or managed grassland (\texttt{E10}, \texttt{E20}).
Organic carbon is used to split points into mineral and organic/peaty
soils using a threshold of 60~g~kg$^{-1}$. The use of
\(\mathrm{pH_{H_2O}}\) as the core spatial discriminator is important:
it provides a direct empirical representation of the continental
acidity gradient, which is the dominant control on lime demand and is
stronger than any purely land-cover-based approximation.

\paragraph{Score definition.}
The point-level lime-demand score is defined by four piecewise-linear
functions of $\mathrm{pH_{H_2O}}$, corresponding to mineral versus
organic/peaty soils and arable versus grassland management. The
functions are derived from RB209 guidance
\citep{RB209,goulding2016}. RB209 is used here as a pragmatic,
transparent proxy basis rather than as a claim that UK agronomic
recommendations apply universally across Europe.

\begin{figure}[htbp]
  \centering
  \includegraphics[width=1.1\textwidth]{figures/Agriculture/RB209.png}
  \caption{RB209-based piecewise lime-demand functions used to define
  the point-level score $s_i^{\mathrm{lime}}$. Separate curves are used
  for mineral and organic/peaty soils, and for arable and grassland
  management.}
  \label{fig:rb209-lime-curves}
\end{figure}

The score is
\begin{equation}
  s_i^{\mathrm{lime}} =
  d_{\mathrm{SOM}(i),\,\mathrm{LU}(i)}\!\left(\mathrm{pH_{H_2O},i}\right),
\end{equation}
where $\mathrm{SOM}(i)\in\{\mathrm{min},\mathrm{org}\}$ and
$\mathrm{LU}(i)\in\{\mathrm{ar},\mathrm{gr}\}$.

\paragraph{Aggregation and normalisation.}
The raw proxy is
\begin{equation}
  \mu_{n,c}^{\mathrm{lime}} =
  \frac{1}{N_{n,c}}
  \sum_{i \in (n,c)} s_i^{\mathrm{lime}},
\end{equation}
with zero assigned when no eligible soil points are available.
Relevance follows Equation~\ref{eq:rho-def}. Relative to the previous
subsections, the main strength of this proxy is that it uses direct soil
measurements rather than an inferred management signal.

\paragraph{Limitations and expected bias.}
RB209 was developed for UK conditions, so its transfer to continental
Europe introduces uncertainty related to mineralogy, texture, and base
saturation. LUCAS Soil 2018 is also a snapshot and cannot capture
long-term re-acidification trends. These limitations are acceptable for
spatial allocation, but they should not be interpreted as a continent-
wide liming recommendation model.

\subsection{Cross-cutting uncertainties and synthesis}

The agricultural proxy block combines datasets with very different
levels of evidential strength. Enteric fermentation, biomass burning,
and soil liming are comparatively well anchored in observable spatial
signals: grazing observations, satellite fire products, and measured
soil pH. By contrast, livestock housing and urea application rely more
heavily on heuristic assumptions because harmonised European data on
housing systems and product-specific fertiliser use are not available at
the required scale.

The main structural uncertainty therefore does not arise from the
mathematical framework itself, but from the degree of interpretation
required to convert available data into source-specific spatial
patterns. The shared framework helps by making those assumptions
explicit: where direct evidence exists, it enters through $\mu$;
where it does not, the heuristic choice is identified and bounded
rather than hidden inside a composite land-cover proxy.

From the perspective of the later case studies, the most important
agricultural upgrades are the NH$_3$-related proxies, because they are
likely to matter most for secondary particulate matter, and the
biomass-burning proxy, because it separates a distinctly episodic source
from the broader agricultural land mask. The CO$_2$ proxies are kept in
the same framework for consistency and completeness, but their main role
here is methodological: they demonstrate that the source-specific
approach can accommodate process controls that differ substantially from
those governing air-pollutant emissions.
